% 컴파일: XeLaTeX 또는 pdfLaTeX 모두 가능 (Overleaf 포함).
% kotex 패키지가 엔진을 자동 감지하여 기본 한글 폰트(나눔/은 글꼴)를 선택합니다.
% Overleaf에서는 좌측 상단 Menu > Compiler 를 XeLaTeX 로 두면 가장 안정적입니다.
\documentclass[10pt]{article}

\usepackage{kotex}

\usepackage[a4paper,margin=2.2cm]{geometry}
\usepackage{enumitem}
\usepackage{xcolor}
\usepackage{listings}
\usepackage{tikz}
\usepackage{hyperref}
\usepackage{titlesec}
\usetikzlibrary{arrows.meta, positioning, fit, backgrounds, calc}

\hypersetup{colorlinks=true, urlcolor=blue!60!black, linkcolor=black}

% 여백/간격 절약
\setlist{nosep, leftmargin=1.4em}
\titlespacing*{\section}{0pt}{1.1ex plus .3ex}{0.6ex}
\titlespacing*{\subsection}{0pt}{0.8ex plus .2ex}{0.4ex}
\setlength{\parskip}{0.35em}
\setlength{\parindent}{0pt}

% 코드 리스팅 스타일 (ASCII 전용)
\definecolor{codebg}{rgb}{0.96,0.96,0.96}
\definecolor{kw}{rgb}{0.10,0.10,0.65}
\definecolor{cmt}{rgb}{0.30,0.55,0.30}
\definecolor{str}{rgb}{0.65,0.20,0.20}
\lstset{
  basicstyle=\ttfamily\scriptsize,
  backgroundcolor=\color{codebg},
  keywordstyle=\color{kw}\bfseries,
  commentstyle=\color{cmt}\itshape,
  stringstyle=\color{str},
  numbers=none, frame=single, framesep=4pt, rulecolor=\color{gray!40},
  breaklines=true, columns=fullflexible, showstringspaces=false,
  aboveskip=0.6em, belowskip=0.4em
}

\title{\vspace{-1.2cm}\textbf{py2cpp: Python\,$\to$\,C++ 정적 변환기 프로젝트 보고서}}
\author{개발자: typeulli \quad / \quad 버전: v1.8.0 \quad / \quad 라이선스: MIT}
\date{2026}

\begin{document}
\maketitle
\vspace{-0.8cm}

%==================================================================
\section{서론 (Introduction)}
%==================================================================
\textbf{py2cpp}는 타입 힌트가 부여된 \emph{제한된 부분집합(restricted subset)}의 Python 코드를 \emph{사람이 읽기 좋은(human-readable)} C++ 소스코드로 변환하는 정적 트랜스파일러(transpiler)이다. 본 프로젝트의 목적은 다음과 같다.

\begin{itemize}
  \item \textbf{인공지능을 사용하지 않는 순수 정적 변환 계층 제공.} LLM 등 확률적 모델이 아닌, AST(추상 구문 트리) 분석과 정적 타입 추론에 기반하여 결정론적이고 재현 가능한 변환을 보장한다.
  \item \textbf{C++ 구조로 깔끔하게 매핑되는 Python 유사 문법 제공.} Python의 타입 힌트를 그대로 활용하여 대응되는 C++ 타입을 생성하며, 가능한 한 많은 Python 문법 요소를 지원한다.
  \item \textbf{Python과 C++ 양쪽 모두의 보일러플레이트(boilerplate) 최소화.} 입력 측에서는 추가 선언을, 출력 측에서는 군더더기 코드를 줄인다.
  \item \textbf{이중 컴파일 모드 제공.} 명령행 도구를 통한 사전 컴파일(AOT, Ahead-of-Time)과 데코레이터를 통한 실시간 컴파일(JIT, Just-in-Time)을 모두 지원하여, 소스 변환과 실행 성능 가속을 동시에 달성한다.
\end{itemize}

즉, py2cpp는 ``Python으로 작성하되 C++의 성능과 가독성을 얻는다''는 목표 아래, 학습용 코드 변환 도구이자 특정 함수의 실행 가속 라이브러리로서 기능하는 것을 목적으로 한다.

%==================================================================
\section{본론 (Main Body)}
%==================================================================

\subsection{전체 프로젝트 아키텍처}
py2cpp는 \emph{(1) 타입 추론 엔진 $\to$ (2) AST 변환 엔진 $\to$ (3) C++ 소스 $\to$ (4) AOT/JIT 소비 경로}의 파이프라인 구조를 가진다. 그림~\ref{fig:arch}는 전체 구조도이다.

\begin{figure}[h]
\centering
\begin{tikzpicture}[
  font=\scriptsize,
  box/.style={draw, rounded corners=2pt, align=center, inner sep=4pt, minimum height=7mm, fill=white},
  stage/.style={draw, rounded corners=3pt, align=center, inner sep=6pt, fill=blue!4, very thick},
  io/.style={draw, align=center, inner sep=4pt, fill=gray!12, minimum height=7mm},
  arr/.style={-{Stealth[length=2mm]}, thick},
  every node/.style={transform shape}
]

% 입력
\node[io] (input) {\textbf{입력: Python 소스 (.py)}\\ 타입 힌트 + py2cpp API\\ \texttt{c\_int, c\_struct, py\_object, @jit} 등};

% Stage 1
\node[stage, below=4mm of input] (s1) {\textbf{(1) 타입 추론 엔진} \;\;\texttt{utils/type.py}\\[2pt]
\begin{tabular}{c}
\texttt{RevealTypeInserter}: 모든 대입/인자에 \texttt{reveal\_type()} 주입 \\
$\Rightarrow$ \texttt{mypy --strict} 실행 $\Rightarrow$ 출력 파싱 $\Rightarrow$ \texttt{TypeContext}
\end{tabular}};

% Stage 2
\node[stage, below=4mm of s1] (s2) {\textbf{(2) AST 변환 엔진} \;\;\texttt{core/compiler.py}\\[2pt]
\begin{tabular}{c}
\texttt{dfs}(문장) / \texttt{dfs\_stmt}(식) 재귀 순회 \;\;|\;\; \texttt{eval\_type}/\texttt{parse\_type} 타입 매핑\\
\texttt{ScriptTemplate}(다중 문장 패턴) \;\;|\;\; \texttt{Builtins}(\#include 의존성 관리)
\end{tabular}};

% C++ 산출물
\node[io, below=4mm of s2] (cpp) {\textbf{(3) C++ 소스코드} (\#include + struct + 전역 + 본문)};

% 소비 경로
\node[box, below left=7mm and -2.2cm of cpp] (aot) {\textbf{(4a) AOT 경로}\\ \texttt{cli.py} (click)\\ $\to$ \texttt{g++} \\ $\to$ \texttt{.cpp / .exe}};
\node[box, below=7mm of cpp] (jit) {\textbf{(4b) JIT 경로}\\ \texttt{core/jit.py} \texttt{@jit}\\ \texttt{g++}$\to$.dll/.so/.dylib \\ \texttt{cl + Python.h}$\to$.pyd\\ \texttt{ctypes} / \texttt{import} + 캐시};
\node[box, below right=7mm and -2.2cm of cpp] (web) {\textbf{(4c) 웹 시뮬레이터}\\ \texttt{api.py} (FastAPI)\\ \texttt{/convert} 엔드포인트\\ pyright LSP (자동완성)};

\draw[arr] (input) -- (s1);
\draw[arr] (s1) -- (s2);
\draw[arr] (s2) -- (cpp);
\draw[arr] (cpp) -- (aot);
\draw[arr] (cpp) -- (jit);
\draw[arr] (cpp) -- (web);

\end{tikzpicture}
\caption{py2cpp 전체 아키텍처 (변환 파이프라인 및 산출물 소비 경로)}
\label{fig:arch}
\end{figure}

\subsection{사용된 소프트웨어 및 하드웨어 기술}
\textbf{개발 언어 및 핵심 의존성}
\begin{itemize}
  \item \textbf{Python 3.10+} — 구조적 패턴 매칭(\texttt{match/case}), \texttt{ParamSpec}, 데이터클래스 등 최신 문법 사용. 표준 라이브러리 \texttt{ast}로 입력 파싱, \texttt{ctypes}로 컴파일 결과 로딩, \texttt{subprocess}로 외부 컴파일러 호출.
  \item \textbf{mypy ($\ge$1.8.0)} — 자체 타입 검사기를 구현하는 대신 mypy를 \emph{타입 추론 백엔드}로 재활용한다. 이것이 본 프로젝트의 핵심 설계 결정이다.
  \item \textbf{click ($\ge$8.1.7)} — 명령행 인터페이스(\texttt{py2cpp} 명령) 구현.
  \item \textbf{FastAPI / uvicorn / pydantic / pyright} — 온라인 시뮬레이터 백엔드(\texttt{api.py}). \texttt{/convert} REST 엔드포인트와, WebSocket 기반 pyright 언어 서버 연동으로 웹 에디터 자동완성을 제공.
\end{itemize}

\textbf{대상(생성) 기술 및 컴파일러}
\begin{itemize}
  \item \textbf{C++17} 표준 코드 생성. \texttt{std::string, std::vector, std::unordered\_map/set, std::tuple, std::complex} 등 STL로 매핑.
  \item \textbf{g++ (MinGW)} — AOT 실행파일 및 JIT용 공유 라이브러리(.dll/.so/.dylib) 빌드.
  \item \textbf{MSVC \texttt{cl.exe} + \texttt{vcvars64}} — CPython 확장 모듈(.pyd) 빌드. \texttt{vswhere.exe}로 Visual Studio 설치 경로를 자동 탐색.
  \item \textbf{CPython C-API (\texttt{Python.h})} — \texttt{py\_object} 타입을 \texttt{PyObject*}로 변환하여 Python 객체를 JIT 함수에 주고받는 상호운용 지원.
\end{itemize}

\textbf{하드웨어 / 실행 환경} — 표준 x86-64 데스크톱 환경(주 개발: Windows). JIT 컴파일 시 최적화 옵션 \texttt{-O3}를 기본 적용하며, \texttt{c\_transform\_par\_unseq}는 \texttt{std::execution::par\_unseq}로 변환되어 멀티코어 병렬 실행을 활용한다.

\subsection{구현 방법 및 알고리즘}

\textbf{(가) mypy 기반 타입 추론.} C++는 정적 타입 언어이므로 모든 변수의 구체 타입이 필요하다. py2cpp는 \texttt{RevealTypeInserter}(\texttt{ast.NodeTransformer})로 모든 대입문·어노테이션·함수 인자·함수 정의 뒤에 \texttt{reveal\_type(x)} 호출을 자동 삽입한 뒤, 이를 임시 파일로 저장하여 \texttt{mypy --strict}를 실행한다. mypy가 출력하는 ``Revealed type is \dots'' 메시지를 정규식으로 파싱하고, AST 경로를 역추적해 변수의 스코프(\texttt{함수명\#변수명})를 식별하여 \texttt{TypeContext}(이름$\to$\texttt{TypeData} 사전)를 구축한다. 즉 \emph{성숙한 타입 검사기를 재활용}하여 추론 정확도를 확보하는 영리한 접근이다.

\textbf{(나) AST 재귀 하강 변환.} 변환은 두 개의 상호 재귀 함수로 이루어진다. \texttt{dfs\_stmt}는 \emph{식}(\texttt{ast.expr})을 C++ 식 문자열로, \texttt{dfs}는 \emph{문장}(\texttt{ast.stmt})을 C++ 문장 문자열로 변환한다. 모듈 최상위에서는 함수 정의를 먼저 끌어올려(hoisting) 출력하고, 나머지 문장은 자동 생성된 \texttt{int main()} 본문에 배치한다. 타입은 \texttt{eval\_type}(국소 추론)과 \texttt{parse\_type}(\texttt{TypeData}$\to$C++ 타입 문자열; \texttt{int$\to$long, float$\to$double, list$\to$std::vector, dict$\to$std::unordered\_map} 등)으로 결정한다.

\textbf{(다) ScriptTemplate 시스템.} ``하나의 Python 식이 여러 C++ 문장을 필요로 하는'' 임피던스 불일치를 해결하는 핵심 기법이다. 예컨대 \texttt{str.split()}, \texttt{list.pop()}, \texttt{sorted()}, 파일 읽기 등은 임시 변수 선언을 동반한다. 각 템플릿은 \texttt{pre}(선행 문장), \texttt{result\_expr}(치환될 식), \texttt{post}(후행 정리 문장)로 구성되며, 충돌 없는 임시 식별자를 난수 생성하여 본문 앞뒤에 코드를 주입한다.

\textbf{(라) Builtins 의존성 관리.} 각 \texttt{Builtins} 노드는 자신이 요구하는 다른 헤더(\texttt{requires})를 알고 있어, 사용된 기능에 따라 필요한 \texttt{\#include}만 위상적으로 수집·중복 제거하여 출력한다. \texttt{namespace} 최소화 옵션(\texttt{using std::vector;} 등)도 함께 처리한다.

\textbf{(마) JIT 컴파일 및 캐싱.} \texttt{@jit} 데코레이터는 \texttt{inspect}로 함수 소스를 추출$\to$데코레이터 제거$\to$C++ 변환$\to$컴파일$\to$동적 로딩한다. 반환/인자가 \texttt{py\_object}이면 \texttt{PyMethodDef/PyModuleDef}를 갖춘 완전한 CPython 확장 모듈(.pyd)을 생성하고, 그렇지 않으면 \texttt{extern "C"} 공유 라이브러리를 생성해 \texttt{ctypes}로 호출한다. \texttt{(py2cpp 버전 + OS + C++ 표준 + 최적화 수준 + 정규화된 AST 소스)}의 MD5 해시를 키로 컴파일 산출물과 메타데이터(JSON)를 OS 캐시 디렉터리에 저장하여, 재실행 시 재컴파일을 생략한다.

\textbf{(바) 지원 범위.} 제어흐름(if/elif/else, for-range, for-each, while, break/continue), 함수, \texttt{@c\_struct} 구조체, 슬라이싱(음수 인덱스 포함), 내장 함수(\texttt{print, input, int, float, abs, len, min, max, sum, sorted, reversed, any, all, round, chr, ord, divmod, open} 등), 문자열/리스트/딕셔너리/집합 메서드, \texttt{math/random/time} 일부 모듈, 파일 입출력(\texttt{with open}), 원시 삽입 지시어(\texttt{c\_include, c\_do, c\_cast, c\_global}), 주석 보존 등을 지원한다.

%==================================================================
\section{수행 결과 (Execution Results)}
%==================================================================
\textbf{AOT 변환 예시.} 다음 Python 코드(좌)는 \texttt{py2cpp input.py} 실행 시 우측의 C++ 코드로 변환된다.

\begin{minipage}[t]{0.49\textwidth}
\begin{lstlisting}[language=Python]
def add(a: int, b: int) -> int:
    return a + b

x = input()
y = 5
print(add(int(x), y))
\end{lstlisting}
\end{minipage}\hfill
\begin{minipage}[t]{0.49\textwidth}
\begin{lstlisting}[language=C++]
#include <iostream>
#include <string>

long add(long a, long b) {
    return a + b;
}
int main() {
    std::string x;
    std::getline(std::cin, x);
    long y = 5;
    std::cout << add(std::stol(x), y)
              << std::endl;
    return 0;
}
\end{lstlisting}
\end{minipage}

타입 힌트(\texttt{int})가 \texttt{long}으로, \texttt{input()}이 \texttt{std::getline}으로, \texttt{int(x)}가 문자열$\to$정수 변환 \texttt{std::stol}로 정확히 매핑되었음을 확인할 수 있다.

\textbf{주요 달성 결과}
\begin{itemize}
  \item \textbf{동작하는 이중 모드 트랜스파일러 완성.} 명령행 AOT(\texttt{.cpp}/\texttt{.exe} 산출)와 데코레이터 기반 JIT가 모두 동작하며, PyPI에 \texttt{pip install py2cpp}로 배포(v1.8.0)되었다.
  \item \textbf{Python C-API 상호운용.} \texttt{py\_object}/\texttt{PyObject*} 변환으로 JIT 함수가 임의의 Python 객체를 인자/반환으로 다룰 수 있다.
  \item \textbf{온라인 시뮬레이터 제공.} \href{https://typeulli.com/py2cpp}{typeulli.com/py2cpp}에서 실시간 변환과 pyright 기반 자동완성을 웹에서 체험 가능.
  \item \textbf{JIT 성능 비교 환경 구축.} 테스트(\texttt{test/jittest.py})에서 1억 회 반복 덧셈 함수를 py2cpp JIT와 Numba로 각각 컴파일·실행하여, 네이티브 C++ 수준의 가속(순수 Python 인터프리터 대비 큰 폭의 실행시간 단축)을 확인하는 벤치마크를 마련하였다.
  \item \textbf{지속적 개발.} v1.2.0부터 v1.8.0까지 점진적으로 지원 문법(문자열/컬렉션 메서드, 파일 I/O, 모듈, 복소수 등)을 확장해 왔다.
\end{itemize}

%==================================================================
\section{결론 (Conclusion)}
%==================================================================
py2cpp는 ``Python의 생산성과 C++의 성능''이라는 상충하는 두 목표를, \emph{AI에 의존하지 않는 결정론적 정적 변환}으로 연결한 프로젝트이다. 특히 자체 타입 검사기를 새로 구현하는 대신 \textbf{mypy를 추론 백엔드로 재활용}하고, 식과 문장의 임피던스 불일치를 \textbf{ScriptTemplate}으로 해소하며, \textbf{Builtins 의존성 그래프}로 헤더를 자동 관리하는 설계는 제한된 자원으로 넓은 문법 범위를 안정적으로 다루기 위한 합리적 선택이었다.

본 프로젝트는 학습용 코드 변환 도구이자 JIT 가속 라이브러리로서 실용적 가치를 가진다. 다만 README가 명시하듯 \emph{성능·메모리 최적화는 수행하지 않으며}, Python 문법의 동적 기능(임의 클래스/상속, 예외 처리, 제너레이터, 동적 타이핑 등)은 지원 범위 밖이라는 한계가 있다. 향후 과제로는 지원 문법 범위 확대, 생성 코드 최적화, 그리고 누락된 일부 런타임 자원(예: \texttt{builtins/} 헬퍼 소스, 웹용 헤더 파일)의 정비를 들 수 있다.

%==================================================================
\section{기타 (Others)}
%==================================================================
\textbf{프로젝트 정보}
\begin{itemize}
  \item \textbf{GitHub 저장소:} \url{https://github.com/typeulli/py2cpp.git}
  \item \textbf{온라인 시뮬레이터:} \url{https://typeulli.com/py2cpp}
  \item \textbf{배포:} PyPI \texttt{py2cpp} (v1.8.0) / 라이선스: MIT (\copyright\ 2025 typeulli) / 규모: 약 4{,}300줄(Python)
\end{itemize}

\textbf{느낀 점 (생각 / 배운 점)}
\begin{itemize}
  \item \textbf{기존 도구의 재활용이 강력하다.} 타입 추론을 처음부터 만드는 대신 mypy의 \texttt{reveal\_type} 출력을 파싱해 활용한 점은, ``바퀴를 재발명하지 않는'' 엔지니어링 판단이 개발 비용과 정확도를 동시에 개선함을 보여준다.
  \item \textbf{언어 간 의미 차이가 변환의 본질적 난점이다.} 단일 Python 식이 여러 C++ 문장으로 풀려야 하는 경우(임시 변수, 정리 코드)를 템플릿의 pre/post 구조로 일반화한 것에서, 트랜스파일러 설계의 핵심이 ``구문 치환''이 아니라 ``의미 보존''에 있음을 배웠다.
  \item \textbf{캐싱과 빌드 환경 추상화의 중요성.} JIT의 해시 기반 캐시와 OS별 컴파일러 탐색(g++/MSVC, \texttt{vswhere})은, 실사용 도구가 되기 위해 ``변환 로직'' 못지않게 ``빌드/배포 환경 처리''가 중요함을 일깨운다.
  \item \textbf{점진적 확장의 가치.} v1.x 동안 문법을 단계적으로 늘려온 이력은, 완벽한 설계를 한 번에 추구하기보다 동작하는 최소 기능에서 출발해 반복적으로 범위를 넓히는 전략의 실효성을 보여준다.
\end{itemize}

\end{document}
